\section{Aritmetika Modular}

Aritmetika modular adalah dasar matematika untuk banyak algoritma kriptografi \cite{trappe2006,menezes1996}. Dua bilangan bulat $a$ dan $b$ dikatakan \textbf{kongruen modulo $n$}, ditulis $a \equiv b \pmod{n}$, jika $n$ membagi $(a - b)$.

\begin{table}[!htbp]
\centering
\caption{Operasi aritmetika modular}
\label{tab:aritmetika-modular}
\begin{tabular}{@{}ll@{}}
\toprule
\textbf{Operasi} & \textbf{Rumus} \\
\midrule
Penjumlahan & $(a + b) \bmod n = ((a \bmod n) + (b \bmod n)) \bmod n$ \\
Pengurangan & $(a - b) \bmod n = ((a \bmod n) - (b \bmod n) + n) \bmod n$ \\
Perkalian & $(a \cdot b) \bmod n = ((a \bmod n) \cdot (b \bmod n)) \bmod n$ \\
\bottomrule
\end{tabular}
\end{table}

\subsection{Invers Modular}
Bilangan $a^{-1}$ adalah invers modular dari $a$ modulo $n$ jika $a \cdot a^{-1} \equiv 1 \pmod{n}$. Invers modular hanya ada jika $\gcd(a, n) = 1$. Algoritma Euclid diperluas digunakan untuk menghitung invers modular \cite{trappe2006,menezes1996}.

\begin{lstlisting}[style=matlab,caption={Operasi modular dalam MATLAB}]
% Contoh: 17 mod 5 = 2
mod(17, 5)   % hasil: 2

% Perkalian modular: (12 * 7) mod 10
mod(12 * 7, 10)   % hasil: 4

% Pangkat modular: 2^10 mod 1000
mod(2^10, 1000)   % hasil: 24
\end{lstlisting}
